\documentclass{article}
\usepackage[T1]{fontenc}
\usepackage{lmodern}
\usepackage{graphicx}
\usepackage{amsmath,amssymb}
\usepackage[a4paper,margin=2cm]{geometry}
\usepackage[absolute,overlay]{textpos}
\setlength{\TPHorizModule}{1mm}
\setlength{\TPVertModule}{1mm}
\usepackage[indonesian]{babel}
\usepackage{hyperref}
\setlength{\emergencystretch}{2em}
% unit: Halaman 30

\begin{document}

% === Halaman 30 (llm) ===

\section*{1. Kode Program RSA (1)}

\begin{verbatim}
from Crypto.PublicKey import RSA
from Crypto.Cipher import PKCS1_OAEP
from Crypto.Util.number import getPrime, inverse

def PembangkitanKunciRSA():
    key_size = 1024    # Gunakan 1024-bit untuk menghindari error
    p = getPrime(key_size // 2)
    q = getPrime(key_size // 2)
    n = p * q
    phi_n = (p - 1) * (q - 1)

    e = 65537    # Kunci publik, default, bisa diganti
    d = inverse(e, phi_n)    # Kunci privat

    # Buat kunci RSA
    key = RSA.construct((n, e, d, p, q))
\end{verbatim}

\end{document}
